\documentclass[11pt,a4paper]{article}
\usepackage[T1]{fontenc}
\usepackage[utf8]{inputenc}
\usepackage[ngerman]{babel}
\usepackage[margin=2.5cm]{geometry}
\usepackage[hidelinks]{hyperref}
\usepackage{array}
\setlength{\parindent}{0pt}
\setlength{\emergencystretch}{2em}
\renewcommand{\arraystretch}{1.25}
\begin{document}
\shorthandoff{"}
\begin{center}
  {\Large\textbf{Webtechnologien -- Übungsblatt 4}}
\end{center}
\section*{Aufgabe 1: Working with Closures}
\textbf{Aufgabe:} Den Code von \texttt{shopA.js} (Warenkorb zu
\texttt{productsA.html}) erklären, zwei Fragen zu Closures und Scopes
beantworten, \texttt{addListener} als anonyme Funktion und als Arrow
Function schreiben, den Button \glqq Empty Shopping Cart\grqq{} umsetzen und
die Produkte nach Preis aufsteigend sortieren. Das HTML darf nicht
geändert werden.

\medskip
\textbf{Erklärung von \texttt{shopA.js}} (Zeilennummern wie in der
Angabe):
\begin{itemize}
  \item Zeile 1: \texttt{'use strict'} schaltet den strengen Modus ein,
    z.\,B. ist eine Zuweisung an eine nicht deklarierte Variable ein Fehler.
  \item Zeile 3: \texttt{products} ist ein Array, in dem für jedes Produkt
    ein Objekt liegt.
  \item Zeile 5--18, \texttt{updateCart()}: leert die Liste
    \texttt{\#orderedItems}, geht alle Produkte durch und erzeugt für jedes
    mit \texttt{numInCart > 0} ein \texttt{li} mit Titel, Preis und Anzahl.
    Dabei wird \texttt{numInCart * price} zur Summe addiert, die am Ende mit
    \texttt{toFixed(2)} in \texttt{\#totalPrice} geschrieben wird.
  \item Zeile 20: Die Schleife geht über alle \texttt{div.product}.
  \item Zeile 21: \texttt{let cp = \{\}} legt pro Durchlauf ein neues, leeres
    Objekt an.
  \item Zeile 23--28: Die Daten werden aus dem HTML gelesen. Das Attribut
    \texttt{data-prodId} wird zu \texttt{dataset.prodid}, weil HTML
    Attributnamen klein schreibt. Für den Titel ist \texttt{firstChild}
    der Textknoten mit dem Zeilenumbruch, \texttt{nextSibling} das
    \texttt{h2}, darin \texttt{firstChild.}\allowbreak\texttt{nodeValue}
    der Text. Den Preis \texttt{"2890,00"} liest \texttt{parseFloat} nur
    bis zum Komma, Ergebnis 2890.
    \texttt{domElement} speichert das \texttt{div} selbst.
  \item Zeile 29: Das Objekt kommt mit \texttt{push} in \texttt{products}.
  \item Zeile 31--35, \texttt{addListener(e)}: gibt das Event aus, erhöht
    \texttt{cp.numInCart} um 1 und ruft \texttt{updateCart()} auf.
  \item Zeile 38--39: Der Listener wird an das \texttt{div.buy} des Produkts
    gehängt. Der Klick auf den Button darin wandert (Bubbling) zum
    \texttt{div} hoch.
\end{itemize}

\textbf{a) a. Warum kennt \texttt{addListener} das richtige
\texttt{cp}?} \texttt{addListener} ist eine Closure: Eine Funktion merkt
sich die Variablen der Umgebung, in der sie \emph{definiert} wurde, nicht
der, in der sie später aufgerufen wird. \texttt{cp} ist mit \texttt{let}
im Schleifenrumpf deklariert und damit block-scoped. Jeder Durchlauf hat
seinen eigenen Block mit einem eigenen \texttt{cp}, und die
\texttt{addListener}-Funktion dieses Durchlaufs hält eine Referenz genau
auf dieses \texttt{cp}. Wenn der Klick viel später kommt, ist die Schleife
längst fertig, die Variable lebt aber weiter, weil die Funktion sie noch
braucht.

\medskip
\textbf{a) b. \texttt{var cp = \{\};} statt \texttt{let cp = \{\};}}
In JavaScript gibt es den globalen Scope, den Funktions-Scope und den
Block-Scope (alles in \texttt{\{ \}}). \texttt{let} und \texttt{const}
gelten nur im Block, \texttt{var} gilt in der ganzen umgebenden Funktion,
hier also im ganzen Script. Mit \texttt{var} gibt es deshalb nur
\emph{eine} Variable \texttt{cp} für alle Durchläufe; jeder Durchlauf
überschreibt sie mit einem neuen Objekt. Die Objekte in \texttt{products}
sind zwar noch verschieden, weil sie vorher mit \texttt{push} gespeichert
wurden, aber alle vier Listener schauen auf dieselbe Variable. Nach der
Schleife zeigt sie auf das letzte Produkt. Ergebnis: Jeder
\glqq add to cart\grqq-Button erhöht die Anzahl von \glqq Tamron
17-24mm\grqq, egal welches Produkt man anklickt (getestet: Klick auf Sony
Alpha IV und Nicon ergibt \glqq Tamron 17-24mm 800(2)\grqq).

\medskip
\textbf{Ergebnis} (\texttt{productsA.html} unverändert): \texttt{shopA.js}
ist die Originaldatei. Was wegfällt, ist nur auskommentiert und mit
\texttt{[1b ENTFERNT]} markiert; alles Neue steht nach fünf Leerzeilen am
Dateiende zwischen \texttt{[NEU]} und \texttt{[ENDE NEU]}.
\begin{itemize}
  \item \textbf{b)} Auskommentiert: \texttt{function addListener} und die
    beiden Zeilen mit \texttt{clickButton} in der Schleife. Neu am Ende: eine Schleife
    \texttt{for (let cp of products)}, die für jedes Produkt den Button aus
    \texttt{cp.domElement} holt und die Funktion direkt in
    \texttt{addEventListener} schreibt. Durch \texttt{let} hat wieder jede
    Runde ihr eigenes \texttt{cp}, die Closure funktioniert wie im Original.
  \item \textbf{b) a.} Standard-Anonyme Funktion (als Kommentar, zum Testen
    einkommentieren und b) b. auskommentieren):
    \begin{center}\texttt{clickButton.addEventListener('click', function (e) \{ ... \})}\end{center}
  \item \textbf{b) b.} Arrow Function (aktiv):
    \begin{center}\texttt{clickButton.addEventListener('click', (e) => \{ ... \})}\end{center}
    Der Inhalt ist in beiden Fällen gleich wie bei \texttt{addListener}.
  \item \textbf{c)} Listener am Button \texttt{\#emptyCard} (so heißt die id
    im HTML): setzt bei allen Produkten \texttt{numInCart = 0} und ruft
    \texttt{updateCart()} auf. Die Liste ist danach leer, Total 0.00.
  \item \textbf{d)} \texttt{products.sort((a, b) => a.price - b.price)}
    sortiert das Array aufsteigend. Danach wird jedes
    \texttt{p.domElement} der Reihe nach mit \texttt{before} vor das
    \glqq Your Order\grqq-\texttt{div} gesetzt. \texttt{before} verschiebt
    ein schon vorhandenes Element, es wird nichts kopiert. Neue Reihenfolge:
    Tamron (800), Nicon (900), Sony 24-105 (1300), Sony Alpha IV (2890).
\end{itemize}

\textbf{Hinweis:} Der Button \glqq add to cart\grqq{} ist ein
\texttt{input}, der Listener hängt aber wie in der Angabe am umgebenden
\texttt{div.buy}; ein Klick daneben im \texttt{div} zählt also auch. Die
Preise werden wegen des Kommas ohne Cents gelesen, das betrifft hier
nichts, weil alle Preise auf \texttt{,00} enden.

\medskip
\textbf{Test:} In Chrome, \texttt{productsA.html} direkt geöffnet:
\begin{center}
\begin{tabular}{|>{\raggedright\arraybackslash}p{5.6cm}|>{\raggedright\arraybackslash}p{8.1cm}|}
\hline
\textbf{Aktion} & \textbf{Ergebnis} \\ \hline
Seite laden & Produkte nach Preis: Tamron, Nicon, Sony 24-105, Sony
  Alpha IV; \glqq Your Order\grqq{} darunter \\ \hline
2x Sony Alpha IV, 1x Nicon & \glqq Nicon DX 16-600 900(1)\grqq,
  \glqq Sony Alpha IV 2890(2)\grqq, Total 6680.00 \\ \hline
Empty Shopping Cart & Liste leer, Total 0.00 \\ \hline
danach 1x Tamron & \glqq Tamron 17-24mm 800(1)\grqq \\ \hline
Variante b) a. aktiv & gleiches Verhalten wie b) b. \\ \hline
Originalcode mit \texttt{var cp} & jeder Button erhöht Tamron \\ \hline
\end{tabular}
\end{center}
Keine Fehler in der Konsole, \texttt{productsA.html} byteweise gleich wie
in der Angabe.

\medskip
{\small \textbf{Hilfsmittel:} Claude Code (Claude Opus 5.5) für Erstellung, Prüfung und Dokumentation.}
\end{document}
